


\section{FPGA to CPU interface Library File - xilinx/cpu.3pl}
\label{chap:lmxilcpulib}

    \index{library!xilinx/cpu.3pl}
    This library provides a basic FPGA-CPU interface for platforms which provide
    a communications path between a CPU and FPGA. This includes the zynq FPGA
    and various FPGA prototype boards which have a USB port for configuration
    loading and serial data communications. These procedures allow the CPU to
    read and write static variables, queues and memories and the FPGA to
    interrupt the CPU.
    
    {\bf This library is included via xilinx/uart\_cpu.3pl
    (\autorefph{chap:lmuartcpulib}) and xilinx/zynq\_axi.3pl
    (\autorefph{chap:lmzynqaxilib}) and should not be included directly by
    application programs.}
    
    xilinx/zynq\_axi.3pl is included in -\\
    boards/avnet/microzed\_7010.3pl\\
    boards/avnet/microzed\_7020.3pl\\
    boards/digilent/zybo.3pl
    
    xilinx/uart\_cpu.3pl is included in -\\
    boards/avnet/aes\_sp3a\_eval400.3pl\\
    board/gadget\_factory/papilio\_one\_100.3pl\\
    board/gadget\_factory/papilio\_one\_250.3pl\\
    board/gadget\_factory/papilio\_one\_500.3pl\\
    board/gadget\_factory/papilio\_pro\_lx4.3pl\\
    board/gadget\_factory/papilio\_pro\_lx9.3pl
    
    This library is incorporated by using a {\bf source} statement -
    \begin{lstlisting}[gobble=4, language=threepl]
       source "xilinx/cpu.3pl"
    \end{lstlisting}
    



    \subsection{library modules}
    


        None.



    \subsection{library classes}
    
        None.


    
    \subsection{library procedures}
            
                


            

        \subsubsection{cpu\_read\_memory - a CPU-readable interface
                                            to a memory}\label{sec:lpcpurdmem}

            {\bf Library: xilinx/cpu.3pl}

            {\bf Prototype:} \\
            \vsp
            \begin{lstlisting}[gobble=12, language=threepl]
               global procedure
               cpu_read_memory (
                   str(name),
                   var(mem),
                   int(port)
               ) {}
            \end{lstlisting}

            {\bf Output parameters:} \\
            none

            {\bf Input parameters:} \\
            {\bf name} is a string identifier for the interface.\\
            {\bf mem} is the cmemory or rmemory variable.\\
            {\bf port} is memory port, 0 or 1, to be used for the read.
        
            {\bf In-line target execution:} \\
            no

            {\bf description:} \\
            This procedure creates a module which allows the CPU to read
            from a registered output memory (rmemory) or a CLB memory (cmemory)
            via the GP1 AXI bus.
            
            The size of the memory data type must not exceed 256 bits as the CPU
            interface software limits burst transfers to 8 32-bit words.
            
            This interface can use any memory port and that port can also be
            used by other application code, but no arbitration of shared access
            via a port is provided. Coincident access will usually result in
            data corruption. If it is not feasible to dedicate a port to CPU
            access then the CPU application code must limit acccess to the
            memory to times when it is known that the FPGA logic will not also
            be accessing the memory port.
            
            On running info2c on the .json file the following C function will have been
            defined for interface {\em name} -
            
            \begin{tabular}{| l | l |}
            \hline
            function                 &                                \\ \hline \hline
            fpga\_{\em name}\_read() & read a datum from FPGA memory  \\ \hline
            \end{tabular}
            
            Example: \\
            If 3PL source file prog.3pl contains the following \\
            \begin{lstlisting}[gobble=12, language=threepl]
               rmemory(m, 2, "uint:8", "uint:16");
               .
               .
               cpu_read_memory("mem1", m, 0);
               .
               .
            \end{lstlisting}
            then for zynq the CPU code should contain something like \\
            \begin{lstlisting}[gobble=12, language=C]
                   #include "prog.h"
                   .
                   .
                   fpga_t *fpga = NULL;
                   int     d;
                   .
                   .
                   nrete( fpga = fpga_open(fpga_unit_default, fpga_flag_default) );
                   erete( fpga_map (fpga) );
                   .
                   .
                   d = mem1_read(fpga, addr);
                   .
            \end{lstlisting}
            where addr is the memory address. The returned data is assigned
            to variable d.

            This procedure creates no in-line target code so the current clock domain
            when this procedure is called is not relevant. The clock domain is not
            changed by this procedure,
            
            {\bf The support CPU software for the UART CPU interface is still
            under development, but application code wil be similar to the
            above.}

            See also - \\
            {\bf cpu\_read\_value()} \autorefph{sec:lpcpurdvalue} \\
            {\bf cpu\_write\_static()} \autorefph{sec:lpcpuwrstat} \\
            {\bf cpu\_read\_queue()} \autorefph{sec:lpcpurdqueue} \\
            {\bf cpu\_write\_queue()} \autorefph{sec:lpcpuwrqueue} \\
            {\bf cpu\_write\_memory()} \autorefph{sec:lpcpuwrmem} \\
            and for zynq -\\
            {\bf dma\_write()} \autorefph{sec:lmzynqwrdma} \\
            {\bf dma\_read()} \autorefph{sec:lmzynqrddma} \\



            

            
                





        \subsubsection{cpu\_read\_queue - a CPU read from a queue
                                        variable}\label{sec:lpcpurdqueue}

            {\bf Library: xilinx/cpu.3pl}

            {\bf Prototype:} \\
            \vsp
            \begin{lstlisting}[gobble=12, language=threepl]
               global procedure
               cpu_read_queue (
                   static(count),
                   queue(out),
                   static(access, "log")
                   <-
                   str(name),
                   queue(in),
                   log(wordcount),
                   var(ilevel)
               ) {}
            \end{lstlisting}

            {\bf Output parameters:} \\
            {\bf count} is an optional queue occupancy count.\\
            {\bf out} is an optional output queue.\\
            {\bf access} is an optional access indicator.\\        

            {\bf Input parameters:} \\            
            {\bf name} is a string identifier for the interface.\\
            {\bf in} is the input queue.\\            
            {\bf wordcount} is an optional flag to enable the number of words
            available in the queue to be determined.\\
            {\bf ilevel} is an optional argument to allow the queue buffer level
            at which an interrupt occurs to be set.\\
        
            {\bf In-line target execution:} \\
            no

            {\bf description:} \\
            This procedure creates a module which is an interface to allow the CPU
            to read a queue variable. The module is created at the end of immediate
            execution from arguments queued by this procedure. The queue can be any
            type, primitive or compound. When the CPU reads from the queue a datum
            is removed from the queue buffer. If the CPU reads from the specified
            address when the queue is empty the read will return all ones .

            The size of the queue data type must not exceed 256 bits as the CPU
            interface software limits burst transfers to 8 32-bit words.

            The queue must be written somewhere in the program (queues without
            sources result in a fatal error on completion of immediate execution).
            When the module is created the argument queue must have an established
            write clock domain or a fatal error is generated.

            If the input queue is read anywhere else in the program (i.e. other than
            in this module) then it must be read in the bus clock (c100) domain
            since that is the clock domain in which it is read in the module created
            by this procedure.

            If the argument to {\bf wordcount} is true the number of words in the
            queue buffer can be interrogated by the CPU. The input queue may have
            any of the allowed buffer sizes. If the queue is asynchronous, i.e. the
            read domain is not the bus clock domain (c100), the count will not be
            exact but will be a guaranteed minimum number of readable data.

            If the argument to {\bf wordcount} is false the data returned from
            reading the number of words in the buffer is 0 or 1 instead of a count.
            A return of 1 means the queue buffer is not empty, i.e. can be read.
            This option is provided to simplify the logic where a simple
            empty/non-empty status is required.

            If {\bf wordcount} has no argument then there will no interface to
            directly read the number of words in the queue buffer, however an
            interrupt may be used instead.            

            Note that an empty queue is guaranteed to return all ones to the CPU
            when read. For most data types this single value will usually be invalid
            and hence in that case it can be used as an indication of an empty
            queue. This can be used instead of explicitly reading the available word
            count.

            Arguments to parameter {\bf ilevel} control provision of a level
            interrupt on queue buffer status as follows -

            If there is no argument then no interrupt is provided.

            If the argument is type "log" with value false then the interrupt will
            be asserted when the queue buffer is not empty.

            If the argument is type "log" with value true then the interrupt will be
            asserted when the number of words in the buffer equals or exceeds a
            level which has been written to the interface from the CPU. The initial
            value of this level is 1. A level value of 0 should not be written to
            the interface as that will result in a continuous level interrupt
            regardless of queue state.

            If the argument is type "uint" or "int" then the interrupt will be
            asserted when the number of words in the buffer equals or exceeds the
            fixed value specified by the integer argument.

            The optional output parameter {\bf count} provides a count of the number
            of data in the queue. The inbuilt function queuewords() can only be used
            in a module which reads the queue, which is usually not the case where
            {\bf cpu\_read\_queue()} is the queue destination. {\bf count} must have
            clock domain {\bf bus\_clk}. If it has no clock domain, {\bf
            cpu\_read\_queue()} will set it to {\bf bus\_clk}. The width of {\bf
            count} must be sufficient to hold the input queue depth.

            If optional output parameter {\bf out} has an argument, each datum read
            from the input queue {\bf in} to the CPU local bus is also written to
            queue {\bf out}. The type of {\bf out} must be such that it can be
            assigned from {\bf in}. {\bf out} should not be allowed to block since
            it cannot stop data being popped from {\bf in} to the CPU bus. While
            {\bf out} in unavailable for writing, no attempt is made to write data
            popped from {\bf in} to it.

            If optional output parameter {\bf access} has an argument, the
            static variable of type "log" is assigned true for one clock cycle
            every time a datum is popped from the input queue {\bf in}. {\bf
            access} may have any clock domain.
            
            On running info2c on the .json file the following C functions will have
            been defined for interface {\em name} -
            
            \begin{tabular}{| l | l | l |}
            \hline
            function                                 & when                   &                         \\ \hline \hline
            fpga\_{\em name}\_read()                 & always                 & read a datum            \\ \hline
            fpga\_{\em name}\_avail\_read()          & wordcount has argument & read number of words    \\ \hline
            fpga\_{\em name}\_avail\_thresh\_write() & ilevel "log" true      & set interrupt threshold \\ \hline
            fpga\_{\em name}\_enable()               & ilevel has argument    & enable interrupt        \\ \hline
            fpga\_{\em name}\_disable()              & ilevel has argument    & disable interrupt       \\ \hline 
            fpga\_{\em name}\_wait()                 & ilevel has argument    & wait on interrupt       \\ \hline  
            \end{tabular}
            
            Example: \\
            If 3PL source file prog.3pl contains the following \\
            \begin{lstlisting}[gobble=12, language=threepl]
               queue(p, "uint:16");
               .
               .
               cpu_read_queue("p1", p, true, false, true);
               .
               .
            \end{lstlisting}
            then for zynq the CPU code should contain something like \\
            \begin{lstlisting}[gobble=12, language=C]
                   #include "prog.h"
                   .
                   .
                   fpga_t *fpga = NULL;
                   int     d[...];
                   int     i = 0;
                   .
                   .
                   nrete( fpga = fpga_open(fpga_unit_default, fpga_flag_default) );
                   erete( fpga_map (fpga) );
                   .
                   .
                   fpga_p1_event_wait(fpga);
                   while (fpga_p1_avail(fpga) != 0)
                       d[i++] = fpga_p1_read(fpga);
                   .
            \end{lstlisting}
            The code will wait until the queue p goes from empty to non-empty
            when an interrupt will release the wait. A loop while the queue
            remains non-empty reads until the queue is empty. Note that since
            this code does not use the actual count of the unread items we could
            have made the echeck argument true in which case the CPU code would
            be identical but logic in the FPGA would have been reduced.
            
            This procedure creates no in-line target code so the current clock
            domain when this procedure is called is not relevant. The clock domain
            is not changed by this procedure,
            
            {\bf The support CPU software for the UART CPU interface is still
            under development, but application code wil be similar to the
            above.}

            See also - \\
            {\bf cpu\_read\_value()} \autorefph{sec:lpcpurdvalue} \\
            {\bf cpu\_write\_static()} \autorefph{sec:lpcpuwrstat} \\
            {\bf cpu\_write\_queue()} \autorefph{sec:lpcpuwrqueue} \\
            {\bf cpu\_read\_memory()} \autorefph{sec:lpcpurdmem} \\
            {\bf cpu\_write\_memory()} \autorefph{sec:lpcpuwrmem} \\
            and for zynq -\\
            {\bf dma\_write()} \autorefph{sec:lmzynqwrdma} \\
            {\bf dma\_read()} \autorefph{sec:lmzynqrddma} \\




        \subsubsection{cpu\_read\_value - CPU read from a
                                                value}\label{sec:lpcpurdvalue}

            {\bf Library: xilinx/cpu.3pl}

            {\bf Prototype:} \\
            \vsp
            \begin{lstlisting}[gobble=12, language=threepl]
               global procedure
               cpu_read_value (
                   static(access, "log")
                   <-
                   str(name),
                   value(in)
               ) {}
            \end{lstlisting}

            {\bf Output parameters:} \\
            {\bf access} is an optional static logical variable which is true
            for one clock cycle following a CPU read.

            {\bf Input parameters:} \\
            {\bf name} is a symbolic identifier string for the bus object.\\
            {\bf in} is a value to be read by the CPU. It takes its type from
            the argument. It must not contain queue reads. It may have any
            clock domain.
        
            {\bf In-line target execution:} \\
            no
            
            {\bf description:} \\
            This procedure creates a module which provides a CPU-readable
            interface to a value of up to 256 bits in size. The
            interface is attached to AXI bus GP0. The module is created
            at the end of immediate execution from arguments queued by this
            procedure. It can be any type, primitive or compound. The optional
            output static logical variable {\bf access} goes true for one clock
            cycle to indicate that a read has occurred.
            
            The size of the input data type must not exceed 256 bits as the CPU
            interface software limits burst transfers to 8 32-bit words.
            
            The clock domain of output {\bf access} need not be the same as
            that of {\bf in}. If {\bf access} has no clock domain it will be
            given the bus clock domain.
            
            The parameters associated with this bus interface, including an
            allocated bus address, are written
            to the {\em design}.info file.
            The utility info2c is used to process this file and create
            an include file {\em design}.h which contains a funcion
            definition to access this interface. for example the 3PL
            code
            \begin{lstlisting}[gobble=12, language=threepl]
               cpu_read_value("alpha", v);
            \end{lstlisting}
            will result in info2c producing a declaration like
            \begin{lstlisting}[gobble=12, language=C]
               uint32_t alpha_read (fpga_t * fpga) {}
            \end{lstlisting}
            
            When the CPU reads a value it is sampling that value, hence data may
            be lost if the sampling rate is lower than the rate at which data
            changes. The data sampling logic ensures that the data is a valid
            sample even when the source is changing at the time of the
            read. This is a primitive interface which depends on additional
            protocols being imposed by the application code. For example to
            guarantee data integrity a write to register the data in a static
            variable prior to reading could be used. For a more complex interface see {\bf cpu\_read\_queue()}
            \autorefph{sec:lpcpurdqueue}.
            
            On running info2c on the .json file the following C function will have been
            defined for interface {\em name} -
            
            \begin{tabular}{| l | l |}
            \hline
            function                   &                 \\ \hline \hline
            fpga\_{\em name}\_read()   & read a datum    \\ \hline
            \end{tabular}
            
            Example: \\
            If 3PL source file prog.3pl contains the following \\
            \begin{lstlisting}[gobble=12, language=threepl]
               static(s, "uint:16");
               .
               .
               cpu_read_value("s1", s);
               .
               .
            \end{lstlisting}
            then for zynq the CPU code should contain something like \\
            \begin{lstlisting}[gobble=12, language=C]
                   #include "prog.h"
                   .
                   .
                   fpga_t *fpga = NULL;
                   int     d;
                   .
                   .
                   nrete( fpga = fpga_open(fpga_unit_default, fpga_flag_default) );
                   erete( fpga_map (fpga) );
                   .
                   .
                   d = fpga_p_read(fpga);
                   .
            \end{lstlisting}
            
            This procedure creates no in-line target code so the current clock domain
            when this procedure is called is not relevant. The clock domain is not
            changed by this procedure,
            
            {\bf The support CPU software for the UART CPU interface is still
            under development, but application code wil be similar to the
            above.}

            See also - \\
            {\bf cpu\_write\_static()} \autorefph{sec:lpcpuwrstat} \\
            {\bf cpu\_read\_queue()} \autorefph{sec:lpcpurdqueue} \\
            {\bf cpu\_write\_queue()} \autorefph{sec:lpcpuwrqueue} \\
            {\bf cpu\_read\_memory()} \autorefph{sec:lpcpurdmem} \\
            {\bf cpu\_write\_memory()} \autorefph{sec:lpcpuwrmem} \\
            and for zynq -\\
            {\bf dma\_write()} \autorefph{sec:lmzynqwrdma} \\
            {\bf dma\_read()} \autorefph{sec:lmzynqrddma} \\
           





        \subsubsection{cpu\_receive\_event - interrupt the CPU}\label{sec:lprcvevent}

            {\bf Library: xilinx/cpu.3pl}

            {\bf Prototype:} \\
            \vsp
            \begin{lstlisting}[gobble=12, language=threepl]
               global procedure
               cpu_receive_event (name, value(sig, "log")) {}
            \end{lstlisting}

            {\bf Output parameters:} \\
            none

            {\bf Input parameters:} \\            
            {\bf name} is the interrupt event identifier string. \\
            {\bf sig} is an interrupt signal
        
            {\bf In-line target execution:} \\
            no

            {\bf description:} \\
            This procedure connects a target variable of type log to the interrupt
            logic to drive the CPU interrupt on execution. The signal may
            be any clock domain and may be true for any period, the false to
            true transition causing the interrupt.
            
            This procedure can be called more than once for the same event number,
            the signals from the instances being ORed. Because the signals are ORed
            a long signal may mask several shorter ones so the interrupt signal
            should be kept short.
            
            On running info2c on the .json file the following C functions will have been
            defined for interface {\em name} -
            
            \begin{tabular}{| l | l |}
            \hline
            function                     &                    \\ \hline \hline
            fpga\_{\em name}\_enable()   & enable interrupt   \\ \hline
            fpga\_{\em name}\_disable()  & disable interrupt  \\ \hline 
            fpga\_{\em name}\_wait()     & wait on interrupt  \\ \hline  
            \end{tabular}
            
            Example: \\
            If 3PL source file prog.3pl contains the following \\
            \begin{lstlisting}[gobble=12, language=threepl]
               value(sig, "log");

               cpu_receive_event("EV1", sig);
               .
               .
               assert(sig <-);
               .
            \end{lstlisting}
            then for zynq the CPU code should contain something like \\
            \begin{lstlisting}[gobble=12, language=C]
                   #include "prog.h"
                   .
                   .
                   fpga_t *fpga = NULL;
                   .
                   .
                   nrete( fpga = fpga_open(fpga_unit_default, fpga_flag_default) );
                   erete( fpga_map (fpga) );
                   .
                   .
                   fpga_EV1_wait(fpga);
                   .
            \end{lstlisting}

            This procedure creates no in-line target code so the current clock domain
            when this procedure is called is not relevant. The clock domain is not
            changed by this procedure,
            
            {\bf The support CPU software for the UART CPU interface is still
            under development, but application code wil be similar to the
            above.}
            
            See also - \\
            {\bf interrupt()} \autorefph{sec:lpinterrupt} \\



            

        \subsubsection{cpu\_write\_memory - a CPU-writable interface
                                            to a memory}\label{sec:lpcpuwrmem}

            {\bf Library: xilinx/cpu.3pl}

            {\bf Prototype:} \\
            \vsp
            \begin{lstlisting}[gobble=12, language=threepl]
               global procedure
               cpu_write_memory (
                   var(mem)
                   <-
                   str(name),
                   int(port)
               ) {}
            \end{lstlisting}

            {\bf Output parameters:} \\
            {\bf mem} is the cmemory or rmemory variable.\\

            {\bf Input parameters:} \\
            {\bf name} is a string identifier for the interface.\\
            {\bf port} is memory port, 0 or 1, to be used for the write. For
            cmemory port 1 is read-only.
        
            {\bf In-line target execution:} \\
            no

            {\bf description:} \\
            This procedure creates a module which allows the CPU to write
            to a registered output memory (rmemory) or a CLB memory (cmemory)
            via the AXI GP1 bus.
            
            The size of the memory data type must not exceed 256 bits as the CPU
            interface software limits burst transfers to 8 32-bit words.
            
            This interface can use any memory port and that port can also be
            used by other application code, but no arbitration of shared access
            via a port is provided. Coincident access will usually result in
            data corruption. If it is not feasible to dedicate a port to CPU
            access then the CPU application code must limit acccess to the
            memory to times when it is known that the FPGA logic will not also
            be accessing the memory port.
            
            On running info2c on the .json file the following C function will have been
            defined for interface {\em name} -
            
            \begin{tabular}{| l | l |}
            \hline
            function                   &                              \\ \hline \hline
            fpga\_{\em name}\_write()  & write a datum to FPGA memory \\ \hline
            \end{tabular}
            
            Example: \\
            If 3PL source file prog.3pl contains the following \\
            \begin{lstlisting}[gobble=12, language=threepl]
               rmemory(m, 2, "uint:8", "uint:16");
               .
               .
               cpu_write_memory(m <- "mem2", 1);
               .
               .
            \end{lstlisting}
            then for zynq the CPU code should contain something like \\
            \begin{lstlisting}[gobble=12, language=C]
                   #include "prog.h"
                   .
                   .
                   fpga_t *fpga = NULL;
                   int     d;
                   .
                   .
                   nrete( fpga = fpga_open(fpga_unit_default, fpga_flag_default) );
                   erete( fpga_map (fpga) );
                   .
                   .
                   fpga_mem2_write(fpga, addr, d);
                   .
            \end{lstlisting}
            where addr is the memory addressd the data to be written.

            This procedure creates no in-line target code so the current clock domain
            when this procedure is called is not relevant. The clock domain is not
            changed by this procedure,
            
            {\bf The support CPU software for the UART CPU interface is still
            under development, but application code wil be similar to the
            above.}

            See also - \\
            {\bf cpu\_read\_value()} \autorefph{sec:lpcpurdvalue} \\
            {\bf cpu\_write\_static()} \autorefph{sec:lpcpuwrstat} \\
            {\bf cpu\_read\_queue()} \autorefph{sec:lpcpurdqueue} \\
            {\bf cpu\_write\_queue()} \autorefph{sec:lpcpuwrqueue} \\
            {\bf cpu\_read\_memory()} \autorefph{sec:lpcpurdmem} \\
            and for zynq -\\
            {\bf dma\_write()} \autorefph{sec:lmzynqwrdma} \\
            {\bf dma\_read()} \autorefph{sec:lmzynqrddma} \\


           



        \subsubsection{cpu\_write\_queue - a CPU write to a queue
                                            variable}\label{sec:lpcpuwrqueue}

            {\bf Library: xilinx/cpu.3pl}

            {\bf Prototype:} \\
            \vsp
            \begin{lstlisting}[gobble=12, language=threepl]
               global procedure
               cpu_write_queue (
                   queue(out)
                   <-
                   str(name),
                   log(freecount),
                   var(ilevel),
                   value(reset, "log")
               ) {}
            \end{lstlisting}

            {\bf Output parameters:} \\
            {\bf out} is the queue written to from the CPU. It takes its type
            from the argument.

            {\bf Input parameters:} \\            
            {\bf name} is a string identifier for the interface.\\
            {\bf freecount} is an optional flag to enable the number of free spaces
            available in the queue to be determined.\\
            {\bf ilevel} is an optional argument to allow the queue buffer level
            at which an interrupt occurs to be set.\\
            {\bf reset} is an optional queue reset.\\
        
            {\bf In-line target execution:} \\
            no

            {\bf description:} \\
            This procedure creates a module which writes to a queue variable from
            the CPU. The module is created at the end of immediate execution from
            arguments queued by this procedure. The queue can be any type, primitive
            or compound (including type null). When the CPU writes to the bus data
            address given by {\bf daddr} a datum is appended to the queue buffer. If
            the CPU writes to the specified address when the queue is full the data
            is lost.

            The size of the queue data type must not exceed 256 bits as the CPU
            interface software limits burst transfers to 8 32-bit words.

            The queue must be read somewhere in the program (queues without
            sinks result in a fatal error on completion of immediate execution).
            When the module is created the argument queue must have an established
            read clock domain or a fatal error is generated.

            If the output queue is written anywhere else in the program (i.e. other
            than in this module) then it must be written in the bus clock domain
            (c100) since that is the clock domain in which it is written in this
            module.

            If the argument to {\bf freecount} is true the number of free entries in
            the queue buffer can be interrogated by the CPU. The output queue may
            have any of the allowed buffer sizes. If the queue is asynchronous, i.e.
            the write domain is not the bus clock domain (c100), the count will not
            be exact but will be a guaranteed minimum number of writable spaces.

            If the argument to {\bf freecount} is false the data returned from
            reading the number of free spaces in the buffer is 0 or 1 instead of a
            count. A return of 1 means the queue buffer is not full, i.e. can be
            written. This option is provided to simplify the logic where a simple
            full/non-full status is required.

            If {\bf freecount} has no argument then there will no interface to
            directly read the space remaining in the queue buffer, however an
            interrupt may be used instead.            

            Arguments to parameter {\bf ilevel} control provision of a level
            interrupt on queue buffer status as follows -

            If there is no argument then no interrupt is provided.

            If the argument is type "log" with value false then the interrupt will
            be asserted when the queue buffer is not full.

            If the argument is type "log" with value true then the interrupt will be
            asserted when the number of free spaces in the buffer equals or exceeds
            a level which has been written to the interface from the CPU. The
            initial value of this level is 1. A level value of 0 should not be
            written to the interface as that will result in a continuous level
            interrupt regardless of queue state.

            If the argument is type "uint" or "int" then the interrupt will be
            asserted when the number of free spaces in the buffer equals or exceeds
            the fixed value specified by the integer argument. Again this should not
            be 0.

            If parameter {\bf reset} provides a means of resetting (clearing) the
            output queue. The argument may have a different clock domain to that of
            the write side of the queue and hence this may be more convenient than
            directly resetting the queue where a change of clock domain may need to
            be explicitly provided.

            Note that a write to a queue variable from {\bf cpu\_write\_queue()} (or
            indeed from anywhere) can be detected within the FPGA by means of the
            inbuilt function {\bf waswritten()} \autorefph{sec:ifwaswritten}. When
            using this function the clock domain must be the same as that of the
            queue write, i.e. {\bf bus\_clk}.
            
            On running info2c on the .json file the following C functions will have
            been defined for interface {\em name} -
            
            \begin{tabular}{| l | l | l |}
            \hline
            function                          & when                   &                            \\ \hline \hline
            fpga\_{\em name}\_write()         & always                 & write a datum              \\ \hline
            fpga\_{\em name}\_avail()         & freecount has argument & read number of spaces free \\ \hline
            fpga\_{\em name}\_avail\_thresh() & ilevel "log" true      & set interrupt threshold    \\ \hline
            fpga\_{\em name}\_enable()        & ilevel has argument    & enable interrupt           \\ \hline
            fpga\_{\em name}\_disable()       & ilevel has argument    & disable interrupt          \\ \hline 
            fpga\_{\em name}\_wait()          & ilevel has argument    & wait on interrupt          \\ \hline  
            \end{tabular}


            
            Example: \\
            If 3PL source file prog.3pl contains the following \\
            \begin{lstlisting}[gobble=12, language=threepl]
               queue(p, "uint:16");
               .
               .
               cpu_write_queue(p <- "p2",, true);
               .
               .
            \end{lstlisting}
            then for zynq the CPU code should contain something like \\
            \begin{lstlisting}[gobble=12, language=C]
                   #include "prog.h"
                   .
                   .
                   fpga_t *fpga = NULL;
                   int     d[...];
                   int     i = 0;
                   int     n = sizeof(d) / sizeof int;
                   .
                   .
                   nrete( fpga = fpga_open(fpga_unit_default, fpga_flag_default) );
                   erete( fpga_map (fpga) );
                   .
                   .
                   fpga_p2_avail_thresh(fpga, n);  // set interrupt level
                   fpga_p2_wait(fpga);             // wait for interrupt
                   for (i=0 ; i<n ; i++) {         // transmit data block
                       fpga_p2_write(fpga, d[i]);  // write item
                   fpga_p2_enable(fpga);           // enable interrupt
                   .
            \end{lstlisting}
            
            or perhaps -
            \begin{lstlisting}[gobble=12, language=threepl]
               queue(p, "uint:16");
               .
               .
               cpu_write_queue(p <- "p2", false);
               .
               .
            \end{lstlisting}
            then for zynq the CPU code should contain something like \\
            \begin{lstlisting}[gobble=12, language=C]
                   #include "prog.h"
                   .
                   .
                   fpga_t *fpga = NULL;
                   int     d[...];
                   int     i = 0;
                   int     n = sizeof(d) / sizeof int;
                   .
                   .
                   nrete( fpga = fpga_open(fpga_unit_default, fpga_flag_default) );
                   erete( fpga_map (fpga) );
                   .
                   .
                   for ( i=0 ; i<n ; i++) {                // transmit data block
                       while (fpga_p2_avail(fpga) == 0)    // poll until not full
                           ;
                       fpga_p2_write(fpga, d[i]);          // write item
                   }
                   .
            \end{lstlisting}

            This procedure creates no in-line target code so the current clock domain
            when this procedure is called is not relevant. The clock domain is not
            changed by this procedure,
            
            {\bf The support CPU software for the UART CPU interface is still
            under development, but application code wil be similar to the
            above.}

            See also - \\
            {\bf cpu\_read\_value()} \autorefph{sec:lpcpurdvalue} \\
            {\bf cpu\_write\_static()} \autorefph{sec:lpcpuwrstat} \\
            {\bf cpu\_read\_queue()} \autorefph{sec:lpcpurdqueue} \\
            {\bf cpu\_read\_memory()} \autorefph{sec:lpcpurdmem} \\
            {\bf cpu\_write\_memory()} \autorefph{sec:lpcpuwrmem} \\
            and for zynq -\\
            {\bf dma\_write()} \autorefph{sec:lmzynqwrdma} \\
            {\bf dma\_read()} \autorefph{sec:lmzynqrddma} \\


            

            
                
        \subsubsection{cpu\_write\_static - a CPU-writable static
                                            variable}\label{sec:lpcpuwrstat}

            {\bf Library: xilinx/cpu.3pl}

            {\bf Prototype:} \\
            \vsp
            \begin{lstlisting}[gobble=12, language=threepl]
               global procedure
               cpu_write_static (
                   static(out), static(access, "log")
                   <-
                   str(name)
               ) {}
            \end{lstlisting}

            {\bf Output parameters:} \\
            {\bf out} is the static variable to be written to by the CPU. It
            takes its type from the argument.The argument may be omitted if only an access signal is
            required.\\
            {\bf access} is a static logical variable which is true for one
            clock cycle following a CPU write to {\bf out}. This argument is
            optional.\\

            {\bf Input parameters:} \\
            {\bf name} is a string identifier for the interface.\\
        
            {\bf In-line target execution:} \\
            no

            {\bf description:} \\
            This procedure creates a module which allows the CPU to write to a
            static variable. The module is created at the end of immediate
            execution from arguments queued by this procedure. The variable can
            be any type, primitive or compound. The associated access static
            logical variable goes true for one clock cycle to indicate that a
            write has occurred. Arguments for any of these output parameters may
            be omitted (but not all).

            The size of the output data type must not exceed 256 bits as the CPU
            interface software limits burst transfers to 8 32-bit words.

            Output parameters {\bf out} and {\bf access} are examined in order
            to determine their clock domains. A parameter not matched by an
            argument is ignored. The clock domain of the first matched parameter
            establishes the default clock domain. If the first matched parameter
            has no clock domain it is given the bus clock domain (c100). Further
            matched output parameters must have the same clock domain or if a
            null clock domain they will be given the default clock domain.

            On running info2c on the .json file the following C function will
            have been defined for interface {\em name} -
            
            \begin{tabular}{| l | l |}
            \hline
            function                   &               \\ \hline \hline
            fpga\_{\em name}\_write()  & write a datum  \\ \hline
            \end{tabular}
            
            Example: \\
            If 3PL source file prog.3pl contains the following \\
            \begin{lstlisting}[gobble=12, language=threepl]
               static(s, "uint:16");
               .
               .
               cpu_write_staticpe(s <- "s2");
               .
               .
            \end{lstlisting}
            then for zynq the CPU code should contain something like \\
            \begin{lstlisting}[gobble=12, language=C]
                   #include "prog.h"
                   .
                   .
                   fpga_t *fpga = NULL;
                   int     d;
                   .
                   .
                   nrete( fpga = fpga_open(fpga_unit_default, fpga_flag_default) );
                   erete( fpga_map (fpga) );
                   .
                   .
                   fpga_s2_write(fpga, d);          // write
                   .
            \end{lstlisting}

            This procedure creates no in-line target code so the current clock domain
            when this procedure is called is not relevant. The clock domain is not
            changed by this procedure,
            
            {\bf The support CPU software for the UART CPU interface is still
            under development, but application code wil be similar to the
            above.}

            See also - \\
            {\bf cpu\_read\_value()} \autorefph{sec:lpcpurdvalue} \\
            {\bf cpu\_read\_queue()} \autorefph{sec:lpcpurdqueue} \\
            {\bf cpu\_write\_queue()} \autorefph{sec:lpcpuwrqueue} \\
            {\bf cpu\_read\_memory()} \autorefph{sec:lpcpurdmem} \\
            {\bf cpu\_write\_memory()} \autorefph{sec:lpcpuwrmem} \\
            and for zynq -\\
            {\bf dma\_write()} \autorefph{sec:lmzynqwrdma} \\
            {\bf dma\_read()} \autorefph{sec:lmzynqrddma} \\





        \subsubsection{interrupt - interrupt the CPU}\label{sec:lpinterrupt}

            {\bf Library: xilinx/cpu.3pl}

            {\bf Prototype:} \\
            \vsp
            \begin{lstlisting}[gobble=12, language=threepl]
               global procedure
               interrupt (var(v)) {}
            \end{lstlisting}

            {\bf Output parameters:} \\
            none

            {\bf Input parameters:} \\            
            {\bf v} is an event interrupt event string identifier or the
            interrupt event number, 0 to 15.
        
            {\bf In-line target execution:} \\
            yes

            {\bf description:} \\
            This procedure
            causes an interrupt event. It executes in one cycle.
            
            The in-line code for this procedure executes in the current clock
            domain which is unchanged on exiting the procedure.
            
            On running info2c on the .json file the following C functions will have been
            defined for interface {\em name} -
            
            \begin{tabular}{| l | l |}
            \hline
            function                    &                   \\ \hline \hline
            fpga\_{\em name}\_enable()  & enable interrupt  \\ \hline
            fpga\_{\em name}\_disable() & disable interrupt \\ \hline 
            fpga\_{\em name}\_wait()    & wait on interrupt \\ \hline  
            \end{tabular}
            
            Example: \\
            If 3PL source file prog.3pl contains the following \\
            \begin{lstlisting}[gobble=12, language=threepl]
               .
               .
               interrupt("EV2");
               .
               .
            \end{lstlisting}
            then \\
            info2c prog \\
            is run, the CPU code should contain something like \\
            \begin{lstlisting}[gobble=12, language=C]
                   #include "prog.h"
                   .
                   .
                   fpga_t *fpga = NULL;
                   .
                   .
                   nrete( fpga = fpga_open(fpga_unit_default, fpga_flag_default) );
                   erete( fpga_map (fpga) );
                   .
                   .
                   fpga_EV2_wait(fpga);
                   .
            \end{lstlisting}
            
            See also \\
            {\bf cpu\_receive\_event()} \autorefph{sec:lprcvevent} \\




    \subsection{library functions}
    
    None.

